\documentclass{cskkuproject}
\usepackage{colortbl} % ใช้ \cellcolor ในตาราง Gantt chart (บทที่ 3)
\usepackage[style=apa,backend=biber]{biblatex} % รูปแบบการอ้างอิง APA 7th edition
\addbibresource{references.bib}

% ==========================================
% ข้อมูลโครงงาน — กรอกข้อมูลที่นี่
% ==========================================

\titleTH{ไดโน : เว็บแอพลิเคชันศูนย์กลางของการท่องเที่ยวจังหวัดขอนแก่น}
\titleEN{Dino: The central web application for tourism in Khon Kaen}

% ข้อมูลนักศึกษาคนที่ 1
\firstStudentID{663380203-9}
\firstStudentPrefixTH{นาย}
\firstStudentPrefixEN{Mr.}
\firstStudentNameTH{คเณศ บัวสด}
\firstStudentNameEN{Khanet Buasod}

% ข้อมูลนักศึกษาคนที่ 2
\secondStudentID{663380228-3}
\secondStudentPrefixTH{นาย}
\secondStudentPrefixEN{Mr.}
\secondStudentNameTH{ภูริ ชาริชา}
\secondStudentNameEN{Phuri Charicha}

% ข้อมูลหลักสูตร
\degreeTH{วิทยาศาสตรบัณฑิต}
\degreeEN{Bachelor of Science}
\majorTH{วิทยาการคอมพิวเตอร์}
\majorEN{Computer Science}

% อาจารย์ที่ปรึกษา
\advisorTH{อ.ธนพล ตั้งชูพงศ์}
\advisorEN{Thanapon Tangchupong}

% ปีการศึกษา
\academicYearTH{2569}
\academicYearEN{2026}

% รหัสวิชา (CP354 771 = โครงงาน 1, CP354 772 = โครงงาน 2)
\courseCode{CP354 771}
\courseName{โครงงานทางวิทยาการคอมพิวเตอร์ 1}
\semester{1}

% คำสำคัญ -- จาก Research Me (1).docx (บทคัดย่อ)
\keywordsTH{แอปพลิเคชันการท่องเที่ยว, ปัญญาประดิษฐ์, การดึงข้อมูลเสริม (RAG), โมเดลภาษาขนาดใหญ่, ไมโครเซอร์วิส, ขอนแก่น}
\keywordsEN{Tourism Application, Artificial Intelligence, Retrieval-Augmented Generation (RAG), Large Language Models (LLMs), Microservices, Khon Kaen}


\begin{document}

% ==========================================
% หน้าปก
% ==========================================
% คำสั่งด้านล่างสร้างหน้าปกอัตโนมัติจากข้อมูลที่กรอกด้านบน
% ไม่ต้องแก้ไขส่วนนี้
\maketitlepage

% ตั้งค่าเลขหน้าส่วนต้นเป็นอักษรไทย (ก, ข, ค, ...)
\pagenumbering{arabic}
\setcounter{page}{0}
\renewcommand{\thepage}{\ThaiAlph{\value{page}}}

% ==========================================
% บทคัดย่อ
% ==========================================
\begin{abstractTH}
% บทคัดย่อภาษาไทย
% ตัวเลขความคืบหน้าปรับปรุงให้ตรงกับบทที่ 5 ฉบับปัจจุบัน (สถานะ ณ เดือนสิงหาคม 2569)
% ส่วนหัว (ชื่อผู้จัดทำ/ปี/ที่ปรึกษา) และคำสำคัญท้ายบทคัดย่อ สร้างอัตโนมัติจาก main.tex ไม่ต้องใส่ในไฟล์นี้

โครงงานนี้มีวัตถุประสงค์เพื่อพัฒนาเว็บแอปพลิเคชัน ``ไดโน'' (Dino) ให้เป็นแพลตฟอร์มศูนย์กลางข้อมูลการท่องเที่ยวในจังหวัดขอนแก่น เพื่อแก้ไขปัญหาการเข้าถึงข้อมูลการท่องเที่ยวที่กระจัดกระจายและไม่ครบถ้วน โครงงานนี้ได้บูรณาการเทคโนโลยีปัญญาประดิษฐ์ (AI) ผ่านโมเดลภาษาขนาดใหญ่ (LLMs) ร่วมกับสถาปัตยกรรม Retrieval-Augmented Generation (RAG) ในการพัฒนาระบบหลัก 2 ส่วน ได้แก่ (1) ระบบสร้างแผนการเดินทางแบบปรับตัว (Adaptive Trip Planning) ที่ประเมินจากข้อจำกัดด้านงบประมาณ ระยะเวลา และความสนใจเฉพาะบุคคล พร้อมกลไกตรวจสอบคุณภาพแผนด้วยเทคนิค LLM-as-a-Judge และ (2) ระบบแชตบอต (Chatbot) สำหรับโต้ตอบและให้คำแนะนำด้านการท่องเที่ยวแบบเรียลไทม์ โดยใช้ฐานข้อมูลสถานที่และร้านอาหารจาก Google Places API ภายในจังหวัดขอนแก่น

ผลการดำเนินงาน ณ ปัจจุบันพบว่า ระบบหลังบ้าน ระบบส่วนหน้า และบริการปัญญาประดิษฐ์เชื่อมต่อและทำงานร่วมกันได้จริงแบบครบวงจร (End-to-End) แล้ว ทั้งฟังก์ชันแชทบอทและระบบวางแผนการเดินทางด้วย LLM ซึ่งมีความคืบหน้ารวมประมาณร้อยละ 85--90 ของขอบเขตที่ออกแบบไว้ อย่างไรก็ตาม ระบบยืนยันตัวตนและระบบสแกน QR สะสม/แลกพอยท์ยังอยู่ในลักษณะต้นแบบ (Prototype) ที่ยังไม่เชื่อมต่อกับฐานข้อมูลจริง และยังไม่ได้ผ่านการทดสอบกับผู้ใช้งานจริงในวงกว้าง (User Acceptance Testing) ทั้งนี้ เมื่อการพัฒนาส่วนที่เหลือและการทดสอบกับผู้ใช้งานจริงเสร็จสมบูรณ์ คาดว่าแอปพลิเคชันนี้จะช่วยลดระยะเวลาในการค้นหาข้อมูล เพิ่มประสิทธิภาพในการวางแผนการเดินทางที่ตอบโจทย์เฉพาะบุคคล ตลอดจนช่วยสนับสนุนผู้ประกอบการท้องถิ่นและส่งเสริมเศรษฐกิจการท่องเที่ยวของจังหวัดขอนแก่นอย่างยั่งยืน

\end{abstractTH}

\begin{abstractEN}
% ABSTRACT (English)
% Progress figures updated to match Chapter 5 as of the current report (status as of August 2026).
% The citation header (authors/year/advisor) and the Keywords line are generated
% automatically from main.tex; do not add them in this file.

This project aims to develop the web application "Dino" as a centralized tourism information platform in Khon Kaen province to address the problem of fragmented and incomplete tourism data. The project integrates artificial intelligence (AI) technologies via Large Language Models (LLMs) combined with Retrieval-Augmented Generation (RAG) architecture to develop two main systems: (1) an Adaptive Trip Planning system that generates itineraries based on constraints regarding budget, time, and personal interests, with plan quality automatically verified through an LLM-as-a-Judge mechanism, and (2) a real-time Chatbot system for interactive communication and tourism recommendations. The system utilizes a database of attractions and restaurants retrieved from the Google Places API within Khon Kaen province.

Current results indicate that the backend, frontend, and AI service are fully connected and operating end-to-end, including the streaming chatbot and the LLM-based trip planner, with an estimated 85--90\% of the designed scope complete. However, user authentication and the QR-based points collection/redemption system remain prototype-level features that are not yet connected to persistent storage, and the system has not yet undergone broad User Acceptance Testing (UAT). Once the remaining components are completed and tested with real users, this application is expected to reduce information search time, enhance the efficiency of personalized trip planning, support local entrepreneurs, and promote the sustainable tourism economy of Khon Kaen province.

\end{abstractEN}

% ==========================================
% กิตติกรรมประกาศ
% ==========================================
\begin{acknowledgement}
% กิตติกรรมประกาศ
% เนื้อหาจาก Research Me (1).docx (กิตติกรรมประกาศ)
% หมายเหตุ: บรรทัดลงชื่อ "ผู้จัดทำ" ต่อท้ายถูกสร้างอัตโนมัติจากชื่อนักศึกษาใน main.tex ไม่ต้องใส่ในไฟล์นี้

ในการดำเนินโครงงานครั้งนี้ ผู้จัดทำโครงงานได้รับความอนุเคราะห์และความช่วยเหลือจากบุคคลหลายท่านด้วยกัน จึงขอขอบคุณไว้ ณ โอกาสนี้ ขอขอบพระคุณอาจารย์ทุกท่านที่ให้ความรู้และอบรมสั่งสอน ขอขอบพระคุณอาจารย์ธนพล ตั้งชูพงศ์ ที่เป็นที่ปรึกษาโครงงานและได้ให้คำชี้แนะแนวทางในการคิดวิเคราะห์และการพัฒนาระบบของโครงงานนี้ ขอขอบคุณคุณพ่อ คุณแม่ ที่เป็นกำลังใจในการเรียนและการทำโครงงานมาโดยตลอด ขอบคุณเพื่อน ๆ ที่ให้คำปรึกษาและช่วยเหลือในด้านการวิเคราะห์ออกแบบระบบ

\end{acknowledgement}

% ==========================================
% สารบัญ
% ==========================================
\tableofcontents
\clearpage
\listoffigures
\clearpage
\listoftables
\clearpage

% ==========================================
% สัญลักษณ์และคำย่อ (ถ้ามี — uncomment ด้านล่าง)
% ==========================================
% \begin{symbolpage}
% \input{symbols}
% \end{symbolpage}

% ==========================================
% เนื้อเรื่อง 5 บท
% ==========================================
% เริ่มนับเลขหน้าใหม่เป็นเลขอารบิก (1, 2, 3, ...)
% ไม่ต้องแก้ไขส่วนนี้
\pagenumbering{arabic}
\setcounter{page}{1}

\chapterone
\chaptertwo
\chapterthree
\chapterfour
\chapterfive

% ==========================================
% เอกสารอ้างอิง
% ==========================================
% *** ไม่ต้องแก้ไขส่วนนี้ ***
% ระบบจะดึงรายการอ้างอิงจากไฟล์ references.bib อัตโนมัติ
% ให้แก้ไขรายการอ้างอิงในไฟล์ references.bib แทน
\clearpage
\makeatletter
\let\ps@plain\ps@empty
\makeatother
\printbibliography[title={เอกสารอ้างอิง},heading=bibintoc]

% ==========================================
% ภาคผนวก
% ==========================================
\input{appendices}

% ==========================================
% ประวัติผู้เขียน
% ==========================================
\begin{biography}
% ประวัติผู้เขียน
% หมายเหตุ: ไม่เกิน 1 หน้ากระดาษ

\section*{นักศึกษาคนที่ 1}

นายคเณศ บัวสด เกิดเมื่อวันที่ 5 พฤศจิกายน 2547
สำเร็จการศึกษาระดับมัธยมศึกษาตอนปลาย จากโรงเรียนเบ็ญจะมะมหาราช เมื่อปีการศึกษา 2565
ปัจจุบันกำลังศึกษาระดับปริญญาตรี หลักสูตรวิทยาศาสตรบัณฑิต
สาขาวิชาวิทยาการคอมพิวเตอร์ วิทยาลัยการคอมพิวเตอร์ มหาวิทยาลัยขอนแก่น เมื่อปีการศึกษา 2566

\section*{นักศึกษาคนที่ 2}

นายภูริ ชาริชา เกิดเมื่อวันที่ 12 เมษายน 2547
สำเร็จการศึกษาระดับมัธยมศึกษาตอนปลาย จากโรงเรียนเบ็ญจะมะมหาราช เมื่อปีการศึกษา 2565
ปัจจุบันกำลังศึกษาระดับปริญญาตรี หลักสูตรวิทยาศาสตรบัณฑิต
สาขาวิชาวิทยาการคอมพิวเตอร์ วิทยาลัยการคอมพิวเตอร์ มหาวิทยาลัยขอนแก่น เมื่อปีการศึกษา 2566

\end{biography}

\end{document}
